\newpage

\section{ numerically stable form of the first-order expression energy factor}\label{sect-numerical-stable-first-order}

we present here an equivalent form of the numerical factor
\begin{align}
  {\cal E}(E_{\gamma_f},E_\gamma) = \frac{\left(e^{-\frac{i}{\hbar}E_{\gamma_f}(t_f-t_i)}-e^{-\frac{i}{\hbar}{ E_\gamma}(t_f-t_i)} \right)}{{E_{\gamma_f}- E_\gamma}}
\end{align}

\subsection{ Notation}

Define

\begin{align}\Delta t=t_f-t_i, \qquad
  \tau=\frac{\Delta t}{\hbar},
\end{align}

and abbreviate

\begin{align}
  {\cal E}(E_{\gamma_f},E_\gamma) = \frac{\left(e^{-\frac{i}{\hbar}E_{\gamma_f}(t_f-t_i)}-e^{-\frac{i}{\hbar}{ E_\gamma}(t_f-t_i)} \right)}{{E_{\gamma_f}- E_\gamma}}=
  \frac{e^{-i\tau E_{\gamma_f}}-e^{-i\tau E_{\gamma}}}{E_{\gamma_f}-E_\gamma}
\end{align}

Throughout, the mathematical $\operatorname{sinc}$ function is defined by

\begin{align}\operatorname{sinc}z=\frac{\sin z}{z}, \qquad
  \operatorname{sinc}(0)=1.
\end{align}

\subsection{ First $\operatorname{sinc}$ transformation}

Use the identity
\begin{align} e^{-ix}-e^{-iy} = -i(x-y)e^{-i(x+y)/2}\operatorname{sinc}\left(\frac{x-y}{2}\right).
\end{align}

It follows that
\begin{importantbox}
  \begin{align} {\cal E}(E_{\gamma_f},E_\gamma)=\frac{e^{-i\tau E_{\gamma_f}}-e^{-i\tau E_\gamma}}{E_{\gamma_f}-E_\gamma} = -i\tau
    e^{-i\tau(E_{\gamma_f}+E_\gamma)/2}        \operatorname{sinc} \left[
      \frac{\tau(E_{\gamma_f}-E_{\gamma})}{2}
    \right].
  \end{align}
\end{importantbox}
\section{ numerically stable form of the second-order expression energy factor}

Consider the expression
\begin{align}
  {\cal F}(E_{\gamma_f}, E_{\gamma_1}, E_\gamma)= \frac{1}{E_{\gamma_1}-E_\gamma} \left[
    \frac{e^{-\frac{i}{\hbar}E_\gamma(t_f-t_i)}
    -e^{-\frac{i}{\hbar}E_{\gamma_f}(t_f-t_i)}}
    {E_{\gamma_f}-E_\gamma}
    -
    \frac{e^{-\frac{i}{\hbar}E_{\gamma_1}(t_f-t_i)}
    -e^{-\frac{i}{\hbar}E_{\gamma_f}(t_f-t_i)}}
    {E_{\gamma_f}-E_{\gamma_1}}
  \right].
\end{align}

Although this expression is finite when two or more of the energies
coincide, its direct numerical evaluation is problematic because the
individual terms contain removable singularities and differences of
nearly equal quantities.

We derive forms based on the sinc function that are better suited for
numerical evaluation.

\subsection{ Notation}

Define

\begin{align}\Delta t=t_f-t_i, \qquad
  \tau=\frac{\Delta t}{\hbar},
\end{align}

and abbreviate

\begin{align}    {\cal F}(E_{\gamma_f}, E_{\gamma_1}, E_\gamma)= \frac{1}{E_{\gamma_1}-E_\gamma} \left[
    \frac{e^{-i\tau E_\gamma}-e^{-i\tau E_{\gamma_f}}}{E_{\gamma_f}-E_\gamma}
    -
    \frac{e^{-i\tau E_{\gamma_1}}-e^{-i\tau E_{\gamma_f}}}{E_{\gamma_f}-E_{\gamma_1}}
  \right].
\end{align}

Throughout, the mathematical $\operatorname{sinc}$ function is defined by

\begin{align}\operatorname{sinc}z=\frac{\sin z}{z}, \qquad
  \operatorname{sinc}(0)=1.
\end{align}

\subsection{ First $\operatorname{sinc}$ transformation}\label{second-regularization}

Use the identity
\begin{align} e^{-ix}-e^{-iy} = -i(x-y)e^{-i(x+y)/2}\operatorname{sinc}\left(\frac{x-y}{2}\right).
\end{align}

It follows that

\begin{align}\frac{e^{-i\tau E_a}-e^{-i\tau E_b}}{E_b-E_a} = i\tau
  e^{-i\tau(E_a+E_b)/2} \operatorname{sinc} \left[
    \frac{\tau(E_a-E_b)}{2}
  \right].
\end{align}

Therefore,
\begin{importantbox}
  \begin{align}
    {\cal F}(E_{\gamma_f},E_{\gamma_1},E_\gamma)
    =
    \frac{i\tau}{E_{\gamma_1}-E_\gamma}
    \left[
      e^{-i\tau(E_\gamma+E_{\gamma_f})/2}
      \operatorname{sinc}
      \left(\frac{\tau(E_\gamma-E_{\gamma_f})}{2}\right)
      -
      e^{-i\tau(E_{\gamma_1}+E_{\gamma_f})/2}
      \operatorname{sinc}
      \left(\frac{\tau(E_{\gamma_1}-E_{\gamma_f})}{2}\right)
    \right].
  \end{align}
  This form explicitly removes the apparent singularities at
  \begin{align} E_{\gamma_f}\to E_\gamma \qquad\text{and}\qquad
    E_{\gamma_f}\to E_{\gamma_1}.
  \end{align}
\end{importantbox}

However, it still contains an apparent difference quotient when

$E_{\gamma_1}\to E_\gamma$.

\subsection{ Limit $E_{\gamma_1}\to E_\gamma$}

Define

\begin{align} g(E, E_{\gamma_f})= e^{-i\tau(E+E_{\gamma_f})/2} \operatorname{sinc}
  \left[
    \frac{\tau(E-E_{\gamma_f})}{2}
  \right].
\end{align}

Then

\begin{align}\mathcal A = i\tau
  \frac{g(E_\gamma, E_{\gamma_f})-g(E_{\gamma_1},E_{\gamma_f})}{E_{\gamma_1}-E_\gamma}.
\end{align}

When $E_{\gamma_1}\to E_\gamma=E$,

\begin{align}\frac{g(E_\gamma,E_{\gamma_f})-g(E_{\gamma_1},E_{\gamma_f})}{E_{\gamma_1}-E_\gamma} \longrightarrow -\frac{\partial g(E_\gamma,E_{\gamma_f})}{\partial E_\gamma}.
\end{align}

Let

\begin{align} z=\frac{\tau(E_\gamma-E_{\gamma_f})}{2}.
\end{align}

Differentiating $g(E_\gamma ,E_{\gamma_f})$ gives

\begin{align} \frac{\partial g(E_\gamma,E_{\gamma_f})}{\partial E_\gamma} = \frac{\tau}{2} e^{-i\tau(E+E_{\gamma_f})/2} \left[
    \operatorname{sinc}'z
    -i\operatorname{sinc}z
  \right].
\end{align}

Hence
\begin{importantbox}
  \begin{align}
    {\cal F} (E_{\gamma_f},E_\gamma,E_{\gamma})
    =
    -\frac{\tau^2}{2}
    e^{-i\tau(E_\gamma+E_{\gamma_f})/2}
    \left[
      \operatorname{sinc}z
      +i\,\operatorname{sinc}'z
    \right]
    ,\quad z= \frac{\tau(E_\gamma-E_{\gamma_f})}{2},\quad \operatorname{sinc}'z = \frac{z\cos z-\sin z}{z^2},\quad \operatorname{sinc}'(0)=0
  \end{align}
  For stable numerical evaluation near the origin, one may use the
  piecewise expression
  \begin{align}
    \operatorname{sinc}'z =
    \begin{cases}
      -\dfrac{z}{3}+\dfrac{z^3}{30}-\dfrac{z^5}{840}
      +\dfrac{z^7}{45360}-\dfrac{z^9}{3991680},
      & |z|<z_{\rm cut},\\[1.2ex]
      \dfrac{z\cos z-\sin z}{z^2},
      & |z|\geq z_{\rm cut}.
    \end{cases}
  \end{align}
  For double-precision calculations, $z_{\rm cut}\sim 10^{-3}$ is a
  conservative choice.  Alternatively, if a numerically stable spherical
  Bessel function is available, one may use
  \begin{align}
    \operatorname{sinc}'z=-j_1(z).
  \end{align}
\end{importantbox}

\subsection {Completely degenerate case}

If
\begin{align} E_\gamma=E_{\gamma_1}=E_{\gamma_f},
\end{align}
then $z=0$, and therefore
\begin{align}\operatorname{sinc}(0)=1, \qquad
  \operatorname{sinc}'(0)=0.
\end{align}

The expression reduces to
\begin{importantbox}
  \begin{align}
    \mathcal F(E_\gamma,E_\gamma,E_\gamma)
    =
    -\frac{\tau^2}{2}e^{-i\tau E_\gamma}
    =
    -\frac{(t_f-t_i)^2}{2\hbar^2}
    e^{-iE_\gamma(t_f-t_i)/\hbar}.
  \end{align}
\end{importantbox}
Thus the apparent singularity when all three energies coincide is
completely removable.

\subsection{Fully regular integral representation}
\begin{importantbox}
  A representation with no energy denominators at all is
  \begin{align}
    {\cal F}(E_{\gamma_f},E_{\gamma_1},E_\gamma)    =
    -\tau^2
    \int_0^1 ds\,
    (1-s)\,
    \exp\left\{
      -i\tau
      \left[
        \frac{1-s}{2}(E_\gamma+E_{\gamma_1})+sE_{\gamma_f}
      \right]
    \right\}
    \operatorname{sinc}
    \left[
      \frac{\tau(1-s)(E_{\gamma_1}-E_\gamma)}{2}
    \right].
  \end{align}
  This identity is exact. It contains no divisions by $ (E_{\gamma_1}-E_\gamma)$, $(E_{\gamma_f}-E_\gamma)$, $(E_{\gamma_f}-E_{\gamma_1}), $
  and all factors remain regular when any or all of these differences
  vanish.
\end{importantbox}
In particular,
\begin{align}\operatorname{sinc}(0)=1,
\end{align}
so no special treatment of a zero argument is required provided that the
numerical implementation of $\operatorname{sinc}$ itself is stable.
